\section{Snapshot score approximation}
\label{sec:Snapshot_score_approximation}
In this section, we prove Lemma~\ref{lem:maxksat-weighted-snapshot-bound}. Given a rounding profile $\Theta= (\ell_0,\bw,\mathcal{I},\br)$, 
we first define several linear constraints on the coefficients $\{c_{\ell }, u_{\ell,j}(\alpha), p_{\ell,jt}(\alpha,\beta)\}$ and combine them to formulate a linear programming, denoted by $\mathsf{LP}(\Theta)$,  that  maximizes the approximation $\rho$.
Finally, an AI agent searches for a good  rounding profile $\Theta$, and exact verification gives $\rho=0.742694813$.


\theoremstyle{plain}
\newtheorem*{restatedsnapshotbound}{Lemma~\ref{lem:maxksat-weighted-snapshot-bound}}
\theoremstyle{definition}

\begin{restatedsnapshotbound}
There exist  a rounding profile $\Theta= (\ell_0,\bw,\mathcal{I},\br)$ and coefficients $ \{c_{\ell
 }, u_{\ell,j}(\alpha), p_{\ell,jt}(\alpha,\beta) \}$ where $3\le \ell\le \ell_0$, $j\in[\ell]$ for $u_{\ell,j}$, $1\le j<t\le\ell$ for $p_{\ell,jt}$, and $\alpha,\beta\in \{+,-\}\times [L]$ so that  for $\rho = 0.742694813$,  we have
  \[
  1-\left(\max_{\alpha}(1-q_\alpha)\right)^{\ell_0+1}\ge\rho,
 \]
 and
 \[
 \rho\cdot \OPT(\Phi_{\le \ell_0})\le  L_{\le\ell_0}^\Theta(\Snap(\Phi_{\le\ell_0})) \le \mathcal{E}^{\Theta}(\Phi_{\le\ell_0}).
 \]
\end{restatedsnapshotbound}

\begin{proof}[Proof of Lemma~\ref{lem:maxksat-weighted-snapshot-bound}]
\leavevmode\par
\paragraph{Step 1. Define the first constraint.}
Fix a rounding profile $\Theta = (\ell_0,\bw,\mathcal{I},\br)$. For $3\le\ell\le\ell_0$, define the score of a clause with labels $\alpha_1,\ldots,\alpha_\ell$ by
\[
  H_\Theta(\alpha_1,\ldots,\alpha_\ell)
  :=c_\ell+\sum_{j=1}^\ell u_{\ell,j}(\alpha_j)
    +\sum_{1\le j<t\le\ell}p_{\ell,jt}(\alpha_j,\alpha_t).
\]
The first constraint requires
\begin{equation}
  H_\Theta(\alpha_1,\ldots,\alpha_\ell)
  \le1-\prod_{j=1}^\ell(1-q_{\alpha_j})
  \label{eq:snapshot-score-upper}
\end{equation}
for every $3\le\ell\le\ell_0$ and every choice of labels.
This constraint ensures that the snapshot score is upper bounded by  the expected number of satisfied clauses under the random rounding. 
Clauses of length one or two already use their exact satisfaction probabilities.


\paragraph{Step 2. Define the second constraint.}
The second constraint is to  guarantee the approximation ratio $\rho$ of  the snapshot score. 



At first, we need some new definitions and notations. 
Fix an optimal assignment $\bx^*$ for $\Phi_{\le\ell_0}$. Let $i(v)$ be the bucket index with bias $\widetilde b_v\in I_{i(v)}$. Group the clauses of $\Phi_{\le\ell_0}$ by their length, literal signs, bucket indices, and variable values under $\bx^*$; and we define the \textit{record} by 
\[
  c=(\ell;s_1,\ldots,s_\ell;i_1,\ldots,i_\ell;\tau_1,\ldots,\tau_\ell),
  \qquad \ell\in[\ell_0],\quad s_j\in\sgnset,\quad i_j\in[L],\quad\tau_j\in\bits.
\]
For each clause $C$, these entries are $s_j=s_j(C)$, $i_j=i(v_j(C))$, and $\tau_j=x^*_{v_j(C)}$. Define the number of clauses with record $c$ by 
$$
n_c := \#\{C\in\Phi_{\le\ell_0}:C\text{ has record } c \text{ under } \mathbf{x}^*\}.
$$  





Write the label $\alpha_j=(s_j,i_j)$ and define the contribution of a record to the snapshot score and the optimal value by
\[
  h_\Theta(c):=\begin{cases}
    q_{\alpha_1},&\ell=1,\\
    1-(1-q_{\alpha_1})(1-q_{\alpha_2}),&\ell=2,\\
    H_\Theta(\alpha_1,\ldots,\alpha_\ell),&3\le\ell\le\ell_0,
  \end{cases}
\]
\[
  o(c):=\1\bigl[\exists j\in[\ell]:\ (s_j,\tau_j)\in\{(+,1),(-,0)\}\bigr].
\]
Therefore, 
summing over all records gives
\begin{equation}
  L_{\le\ell_0}^\Theta(\Snap(\Phi_{\le\ell_0}))=\sum_c n_ch_\Theta(c),
  \qquad \OPT(\Phi_{\le\ell_0})=\sum_c n_co(c).
  \label{eq:snapshot-record-sums}
\end{equation}

The bucket bounds restrict the counts $n_c$.
Let $a_i,b_i$ be the endpoints of $I_i$.
Put $\chi(+)=1$ and $\chi(-)=-1$. For $i\in[L]$ and $\tau\in\bits$, define the contribution of a record to the lower and upper bucket bounds by
\[
  g^{\rm lo}_{i,\tau}(c):=
  w_\ell\sum_{\substack{j\in[\ell]:\,\\
  i_j=i,  \tau_j=\tau}}(a_i-\chi(s_j)),
  \qquad
  g^{\rm up}_{i,\tau}(c):=
  w_\ell\sum_{\substack{j\in[\ell]:\\ i_j=i,  \tau_j=\tau } }(\chi(s_j)-b_i),
\]
where $w_\ell$, the $\ell$-th entry of $\bw$, is the weight of clauses with length $\ell$.
Their sums over the records are nonpositive, as the following claim shows.
\begin{claim}\label{clm:snapshot-bucket-sums}
For every $i\in[L]$ and $\tau\in\bits$,
\[
  \sum_c n_cg^{\rm lo}_{i,\tau}(c)\le0,
  \qquad \sum_c n_cg^{\rm up}_{i,\tau}(c)\le0.
\]
\end{claim}
\begin{proof}
Within the variables with $i(v)=i$ and $x_v^*=\tau$, let
\[
  d_{i,\tau}:=\sum_{\substack{v:\,i(v)=i, x_v^*=\tau}}(\widetilde p_v+\widetilde q_v),
  \qquad
  e_{i,\tau}:=\sum_{\substack{v:\,i(v)=i, x_v^*=\tau}}(\widetilde p_v-\widetilde q_v).
\]
Since $a_i\le\widetilde b_v\le b_i$ and $\widetilde p_v-\widetilde q_v=\widetilde b_v(\widetilde p_v+\widetilde q_v)$,
$
  a_id_{i,\tau}\le e_{i,\tau}\le b_id_{i,\tau};
$
each literal in a clause of length $\ell$ contributes $w_\ell$ to $d_{i_j,\tau_j}$ and $w_\ell\chi(s_j)$ to $e_{i_j,\tau_j}$. Hence
\[
  \sum_c n_cg^{\rm lo}_{i,\tau}(c)=a_id_{i,\tau}-e_{i,\tau}\le0,
  \qquad
  \sum_c n_cg^{\rm up}_{i,\tau}(c)=e_{i,\tau}-b_id_{i,\tau}\le0. \qedhere
\]
\end{proof}




\medskip
The second constraint is to   find    coefficients $\lambda^{\rm lo}_{i,\tau}\ge 0,\lambda^{\rm up}_{i,\tau}\ge 0$ such that for   every record $c$,
\begin{equation}
  h_\Theta(c)-\rho\cdot  o(c)
  +\sum_{i,\tau}\bigl(
    \lambda^{\rm lo}_{i,\tau}g^{\rm lo}_{i,\tau}(c)
    +\lambda^{\rm up}_{i,\tau}g^{\rm up}_{i,\tau}(c)\bigr)\ge0.
  \label{eq:snapshot-record-lower}
\end{equation}
This constraint  builds on the bucket bounds and allows more flexibility in the coefficients so that we can achieve a better approximation ratio $\rho$. Without this  relaxation,  if  we simply ask for   
\[
 h_\Theta(c)-\rho\cdot  o(c)\ge 0,
\] the satisfied unary records $(1;\,+;\,i;\,1)$ and $(1;\,-;\,i;\,0)$ would require $r_i\ge\rho$ and $1-r_i\ge\rho$, forcing $\rho\le1/2$.

\paragraph{Step 3. Maximize the approximation ratio.}
Given the rounding profile  $\Theta$, we define a linear program $\mathsf{LP}(\Theta)$, combining the constraints from the last two steps, to maximize the approximation ratio $\rho$; the variables of $\mathsf{LP}(\Theta)$ are $\rho$, the score coefficients $c_\ell,u_{\ell,j}(\alpha),p_{\ell,jt}(\alpha,\beta)$, and the extra  coefficients $\lambda^{\rm lo}_{i,\tau},\lambda^{\rm up}_{i,\tau}$ from Step 2.  
$\mathsf{LP}(\Theta)$ is defined as follows: 
\[
\begin{aligned}
  \text{maximize}\quad &\rho\\
  \text{s.t.}\quad
  & (1).\;H_\Theta(\alpha_1,\ldots,\alpha_\ell)
    \le1-\prod_{j=1}^\ell(1-q_{\alpha_j})
    &&(3\le\ell\le\ell_0;\text{ all labels}),\\
  &(2).\; h_\Theta(c)-\rho o(c)
    +\sum_{i,\tau}\bigl(\lambda^{\rm lo}_{i,\tau}g^{\rm lo}_{i,\tau}(c)
    +\lambda^{\rm up}_{i,\tau}g^{\rm up}_{i,\tau}(c)\bigr)\ge0
    &&\text{for every record }c,\\
  &(3).\; \lambda^{\rm lo}_{i,\tau},\lambda^{\rm up}_{i,\tau}\ge0
    &&(i\in[L],\ \tau\in\bits),\\
  &(4).\; 0\le\rho\le 1-\left(\max_\alpha(1-q_\alpha)\right)^{\ell_0+1}.
\end{aligned}
\]
Once $\Theta$ is fixed, all $q_\alpha$ and $g^{\rm lo}_{i,\tau}(c),g^{\rm up}_{i,\tau}(c)$ are constants. The first two constraint families are therefore linear, and their index sets are finite. The constraint $(4)$ ensures that every clause of length greater than $\ell_0$ is satisfied with probability at least $\rho$ under this rounding.

\paragraph{Step 4. Prove the approximation guarantee.}
We first  prove the upper bound. Take any feasible solution of $\mathsf{LP}(\Theta)$. 
For each $3\le\ell\le\ell_0$, grouping clauses by their labels gives
\begin{equation}
  \label{eq:the_2nd_part_score}
  \begin{aligned}
  \sum_{C\in\Phi: |C|=\ell} H_\Theta(\alpha_1(C),\ldots,\alpha_\ell(C))
  =c_\ell M_\ell+\!\sum_{j=1}^\ell\sum_\alpha u_{\ell,j}(\alpha)U_{\ell,j,\alpha}
  +\sum_{1\le j<t\le\ell}\sum_{\alpha,\beta}
    p_{\ell,jt}(\alpha,\beta)P_{\ell,j,t,\alpha,\beta}.
\end{aligned}
\end{equation}
Under the  random  rounding by $\Theta$, the expected number of satisfied clauses is
\[
\begin{aligned}
  \mathcal E^\Theta(\Phi_{\le\ell_0})
  &=\sum_{\ell=1}^{\ell_0}\sum_{C\in\Phi:\,|C|=\ell}\E[C(\bx)]
  =\sum_{\ell=1}^{\ell_0}\sum_{C\in\Phi:\,|C|=\ell}
    \left(1-\prod_{j=1}^\ell(1-q_{\alpha_j(C)})\right)\\
  &\ge \mathcal{E}^{\Theta}(\Phi_{\le 2})
    +\sum_{\ell=3}^{\ell_0}\sum_{C\in\Phi:\,|C|=\ell}
      H_\Theta(\alpha_1(C),\ldots,\alpha_\ell(C))\\
  &=L_{\le\ell_0}^\Theta(\Snap(\Phi_{\le\ell_0})).
\end{aligned}
\]
The second last inequality is by the constraint  (1), and the last equality is by Equation~\eqref{eq:expect_2_term} and Equality~\eqref{eq:the_2nd_part_score}.





For the lower bound, multiply Inequality~\eqref{eq:snapshot-record-lower} by $n_c$ and sum over records $c$; by  Equality~\eqref{eq:snapshot-record-sums}, we have
\[
\begin{aligned}
  L_{\le\ell_0}^\Theta(\Snap(\Phi_{\le\ell_0}))-\rho\,\OPT(\Phi_{\le\ell_0})+ \sum_{i,\tau} \underbrace{  \left(
    \lambda^{\rm lo}_{i,\tau}\sum_c n_cg^{\rm lo}_{i,\tau}(c)
    +\lambda^{\rm up}_{i,\tau}\sum_c n_cg^{\rm up}_{i,\tau}(c)\right)}_{\le 0}  \ge 0.
\end{aligned}
\]
where the underbrace bound $\le 0$ follows from the nonpositivity of the sums in Claim~\ref{clm:snapshot-bucket-sums} and $\lambda^{\rm lo}_{i,\tau},\lambda^{\rm up}_{i,\tau}\ge 0$. This proves \[L_{\le\ell_0}^\Theta(\Snap(\Phi_{\le\ell_0}))-\rho\,\OPT(\Phi_{\le\ell_0})\ge 0.\]


\paragraph{Step 5. Select and verify the rounding profile.}
An AI agent alternates between updating $\Theta$ and solving $\mathsf{LP}(\Theta)$ to improve $\rho$.



We convert the final coefficients to rational numbers and verify every constraint using exact arithmetic. The resulting rounding profile has $\ell_0=5$, $L=1000$, and for any $i\in [1000]$,  $61/250\le r_i\le189/250$. It gives a feasible solution with
\[
  \rho=0.742694813,
  \qquad 1-\left(\max_\alpha(1-q_\alpha)\right)^{\ell_0+1} =  1-\left(\frac{189}{250}\right)^6>
  \rho.
\]
The code for selecting and verifying this rounding profile is available at \url{https://github.com/Guangxu-Yang/maxksat-lp-certificate}.
\end{proof}
